% =====================================================================
% PARTE I — Contexto, problema e objetivos      [~8 min]
% Fonte: Chapters/introducao, Chapters/contexto_cern_atlas
% =====================================================================

\section{Contexto e Problema}

\begin{frame}{O experimento ATLAS e o calorímetro LAr}
  \duascolunas{%
    \begin{itemize}
      \item Colisões a \dest{40\,MHz}, vazão da ordem de \dest{70\,TB/s}.
      \item O calorímetro de argônio líquido (\LAr) fornece as grandezas
            usadas na reconstrução de energia.
      \item O monitoramento de qualidade de dados (\DQM) decide o que é
            aproveitável para física.
    \end{itemize}
    \pendente{escolher a figura do detector --- ver figures/atlas/}
  }{%
    \centering
    \pendente{\texttt{\textbackslash includegraphics} da figura do LAr}
  }
\end{frame}

\begin{frame}{O gargalo do monitoramento de qualidade de dados}
  \begin{itemize}
    \item Inspeção majoritariamente \dest{manual} e por limiares fixos.
    \item O critério Toast e suas limitações.
      \pendente{enunciar as três limitações do Cap. 3}
    \item Defeitos D3--D11: heterogêneos em escala, duração e assinatura.
    \item Consequência: latência de diagnóstico e perda de dados válidos.
  \end{itemize}
\end{frame}

\begin{frame}{Questão de pesquisa e hipótese}
  \begin{block}{Questão}
    \pendente{transcrever a questão de pesquisa do Cap. 1}
  \end{block}
  \vspace{0.5em}
  \begin{alertblock}{Hipótese}
    \pendente{transcrever a hipótese central --- a sinergia entre detecção,
    reconstrução e previsão supera os agentes isolados}
  \end{alertblock}
\end{frame}

\begin{frame}{Objetivos}
  \textbf{\color{peedark}Objetivo geral}
  \begin{itemize}
    \item \pendente{objetivo geral do Cap. 1}
  \end{itemize}
  \vspace{0.6em}
  \textbf{\color{peedark}Objetivos específicos}
  \begin{enumerate}
    \item \pendente{OE1}
    \item \pendente{OE2}
    \item \pendente{OE3}
  \end{enumerate}
\end{frame}
